\textbf{Predstavil: }

\begin{enumerate}
  \item Za $n \in \mathbb{N}$ in $x \in \mathbb{R}$, naj bo \[f_{n}(x) = \frac{x}{1+nx^{2}}.\]  Dokažite, da zaporedja funkcij $\left\{ f_{n} \right\}$ enakomerno konvergira proti funkciji $f$, in da \[f'(x) = \lim_{n \to \infty} f_{n}'(x)\] velja, če je $x \neq 0$, pa ne, če je $x=0$.  
  \item Obstaja izrek, ki pravi, kdaj velja enačba \[f'(x) = \lim_{n \to \infty} f_{n}'(x)\] za zaporedje funkcij $\left\{ f_{n} \right\}$. Zakaj ta izrek ne drži v prejšnjem primeru?
\end{enumerate}